% ==================== 7.3 研究展望 ====================
\section{研究展望}
\label{sec:prospects}

本文在月面巡视器自主行走数字孪生技术方面取得了一定进展，但仍有许多问题值得进一步研究。未来的研究方向包括：

\textbf{（1）数字孪生模型的精度提升}

本文建立的数字孪生模型在精度方面仍有提升空间。未来可以：
\begin{itemize}
    \item 引入更精细的环境建模方法，如基于物理的月壤生成模型、考虑月壤-车轮耦合的细观力学模型等；
    \item 发展数据驱动的模型更新方法，利用在轨实测数据持续优化模型参数；
    \item 研究模型不确定性的量化与传播方法，提高预测的可靠性。
\end{itemize}

\textbf{（2）多巡视器协同数字孪生}

本文主要研究单巡视器的数字孪生，未来可以扩展到多巡视器协同场景：
\begin{itemize}
    \item 研究多巡视器协同探测的任务分配与调度方法；
    \item 开发多巡视器协同建图与定位算法；
    \item 构建多巡视器协同数字孪生平台，实现协同任务仿真与决策支持。
\end{itemize}

\textbf{（3）人在环路的数字孪生系统}

本文主要关注巡视器的自主决策，未来可以引入人在环路机制：
\begin{itemize}
    \item 研究人机协同决策方法，实现人与数字孪生系统的有效交互；
    \item 开发沉浸式可视化界面，提高地面控制人员的态势感知能力；
    \item 设计人在环路的仿真验证方法，优化任务规划和应急响应策略。
\end{itemize}

\textbf{（4）智能化程度提升}

本文的智能算法主要针对路径规划，未来可以向更多领域拓展：
\begin{itemize}
    \item 研究基于深度学习的环境感知方法，提高岩石识别、地形分类的准确性；
    \item 开发智能故障诊断与健康管理方法，提高巡视器的可靠性和可维护性；
    \item 探索知识图谱、大语言模型等技术在数字孪生中的应用，提升系统的智能化水平。
\end{itemize}

\textbf{（5）星载边缘计算}

本文的数字孪生计算主要在地面进行，未来可以向星载边缘计算发展：
\begin{itemize}
    \item 研究轻量化数字孪生模型，降低对计算资源的需求；
    \item 开发星载边缘计算架构，实现部分数字孪生功能在巡视器端运行；
    \item 研究地-星协同计算方法，实现计算任务的合理分配。
\end{itemize}

\textbf{（6）其他天体探测应用}

本文的方法可推广至火星、小行星等其他天体探测任务：
\begin{itemize}
    \item 研究火星环境数字孪生建模方法，考虑火星大气、沙尘暴等特殊因素；
    \item 开发小行星着陆数字孪生系统，考虑微重力、不规则形状等挑战；
    \item 建立通用的深空探测器数字孪生平台，支撑多种探测任务。
\end{itemize}

总之，月面巡视器自主行走数字孪生技术是一个具有广阔应用前景的研究方向。随着航天技术、人工智能技术和数字孪生技术的不断发展，该领域将迎来更多突破，为人类深空探测事业做出更大贡献。
